\documentclass[11pt]{article}

\usepackage[a4paper,margin=2.2cm]{geometry}
\usepackage{amsmath,amssymb}
\usepackage{siunitx}
\usepackage{microtype}
\usepackage{booktabs}
\usepackage{hyperref}

\title{Physical reading of the exponential recession-tail fit\\
       \large Linearised Richards equation and the swale dataset}
\author{Alexis Shakas \\ \small{ETH Z\"urich, shakasa@ethz.ch}}
\date{\today}

\newcommand{\dpartial}[2]{\frac{\partial #1}{\partial #2}}

\begin{document}

\maketitle

\section{Richards' equation in one dimension}

For a homogeneous, isothermal soil column with vertical coordinate $z$
(positive downward), volumetric water content $\theta(z,t)$,
unsaturated hydraulic conductivity $K(\theta)$ and matric head
$h(\theta)$, mass conservation combined with the Buckingham--Darcy law
gives the mixed-form Richards equation
\begin{equation}
    \dpartial{\theta}{t}
    \;=\;
    \dpartial{}{z}\!\left[
        D(\theta)\,\dpartial{\theta}{z} \;-\; K(\theta)
    \right]
    \;-\;
    S(\theta,z,t),
    \label{eq:richards}
\end{equation}
where $D(\theta) = K(\theta)\,\partial h/\partial\theta$ is the
hydraulic diffusivity and $S(\theta,z,t)$ is the root-uptake / direct
evaporation sink. The nonlinearity is concentrated in $D(\theta)$ and
$K(\theta)$, both of which vary by several orders of magnitude across
the unsaturated range.

\section{Linearisation around a wetting peak}

After a rainfall event the soil moisture relaxes toward an
approximately uniform baseline $\theta_b(z)$. Write
\begin{equation}
    \theta(z,t) \;=\; \theta_b(z) + \theta'(z,t),
    \qquad
    |\theta'| \ll \theta_b.
\end{equation}
Freezing the soil-hydraulic functions at $\theta_b$ and dropping terms
of order $\theta'^2$ converts Eq.~\eqref{eq:richards} into a
linear advection--diffusion--reaction equation in $\theta'$:
\begin{equation}
    \dpartial{\theta'}{t}
    \;=\;
    D_b\,\dpartial{^2\theta'}{z^2}
    \;-\;
    K_b'\,\dpartial{\theta'}{z}
    \;-\;
    \lambda\,\theta',
    \label{eq:linrichards}
\end{equation}
with constant coefficients
\begin{equation}
    D_b = D(\theta_b),
    \qquad
    K_b' = \frac{dK}{d\theta}\bigg|_{\theta_b},
    \qquad
    \lambda = \frac{\partial S}{\partial\theta}\bigg|_{\theta_b} \;\ge\; 0.
\end{equation}

\section{Eigenvalue spectrum and the single-mode regime}

Substituting the standard advection-removing transform
$\theta'(z,t) = \exp\!\big(\frac{K_b' z}{2 D_b}\big)\,\psi(z,t)$ converts
Eq.~\eqref{eq:linrichards} into a pure diffusion equation with a
constant-rate sink:
\begin{equation}
    \dpartial{\psi}{t}
    \;=\;
    D_b\,\dpartial{^2\psi}{z^2}
    \;-\;
    \left[\,\frac{K_b'^{\,2}}{4 D_b} + \lambda\,\right]\!\psi.
\end{equation}
On a layer of effective depth $L$, the eigenvalues of the slowest
decaying spatial mode (the value of which depends on the boundary
conditions chosen at $z=0$ and $z=L$) take the canonical form
\begin{equation}
    \boxed{\;
    \frac{1}{\tau_n}
    \;=\;
    \underbrace{D_b\!\left(\frac{n\pi}{L}\right)^{\!2}}_{\text{vertical
        diffusion mode } n}
    \;+\;
    \underbrace{\frac{K_b'^{\,2}}{4 D_b}}_{\text{gravity
        advection}}
    \;+\;
    \underbrace{\vphantom{\frac{K_b'^{\,2}}{4 D_b}}\lambda}_{\text{ET sink}},
    \;}
    \label{eq:tau-sum}
\end{equation}
where $n=1,2,3,\dots$ The recession is dominated by the fundamental
mode ($n=1$) within a few $\tau_1$ of the peak. With
mixed Dirichlet--Neumann boundaries (closed surface, fixed-$\theta$
bottom into a water table) the prefactor on the diffusion term reads
$\pi^2 D_b/(4 L^2)$ instead of $\pi^2 D_b/L^2$. Either way, the
\emph{form} of Eq.~\eqref{eq:tau-sum} is the same: a single
characteristic timescale built additively from three physical
processes.

\paragraph{Diagnostic.} If a recession is well described by a single
$\tau$ (high R$^2$ for $\theta'(t) = A\,e^{-t/\tau} + C$) and the
corresponding power-law exponent $\alpha$ is near zero, the soil is in
this single-eigenmode regime. Higher modes have already decayed and
no heterogeneity-induced heavy tail is visible.

\section{Comparison with alternative tail models}

Different physical regimes produce qualitatively different recession
shapes:

\begin{center}
\renewcommand{\arraystretch}{1.25}
\begin{tabular}{@{} l l l @{}}
\toprule
Regime & Recession shape & Diagnostic exponent \\
\midrule
Linearised Richards (single mode) & $\theta \propto e^{-t/\tau}$ & $\alpha \approx 0$ \\
Capillary diffusion, semi-infinite half-space & $\theta \propto t^{-1/2}$ & $\alpha \approx 0.5$ \\
Boussinesq groundwater, late time & $\theta \propto t^{-2}$ & $\alpha \approx 2$ \\
Brutsaert--Nieber, $dQ/dt = -aQ^b$ & $\theta \propto t^{-1/(b-1)}$ & $\alpha = 1/(b-1)$ \\
Distribution of pore-drainage times & $\sum_i A_i e^{-t/\tau_i}$ & apparent power-law over a window \\
\bottomrule
\end{tabular}
\end{center}

For the swale data the median fits give $\alpha = 0.006$--$0.04$ and
exponential R$^2$ consistently above the power-law R$^2$ across all
(depth, treatment) cells. This places the system in the single-mode
linear-reservoir regime and justifies using $\tau$ as the primary
descriptor.

\section{Inverting the fitted \texorpdfstring{$\tau$}{tau} for soil parameters}

When the three contributions to Eq.~\eqref{eq:tau-sum} cannot be
separated from a single trace, the fitted $\tau$ provides only an
\emph{upper bound} on the diffusion timescale:
\begin{equation}
    D_b \;\le\; \frac{L^2}{\pi^2\,\tau}.
    \label{eq:D-bound}
\end{equation}
For $L = \SI{10}{cm}$ and the swale-dataset medians:

\begin{center}
\begin{tabular}{@{} l l c r r @{}}
\toprule
Depth & Treatment & $\tau$ (h) & $\tau$ (s) & $D_b$ upper bound (\si{\meter\squared\per\second}) \\
\midrule
\SI{10}{cm} & control & 56   & $2.0\times 10^{5}$ & $5.0\times 10^{-9}$ \\
\SI{10}{cm} & swale   & 85   & $3.1\times 10^{5}$ & $3.3\times 10^{-9}$ \\
\SI{40}{cm} & control & 73   & $2.6\times 10^{5}$ & $3.9\times 10^{-9}$ \\
\SI{40}{cm} & swale   & 3261 & $1.2\times 10^{7}$ & $8.6\times 10^{-11}$ \\
\bottomrule
\end{tabular}
\end{center}

\noindent The control 10 / 40 \si{cm} numbers ($\sim 10^{-9}\,\si{m^2/s}$) sit in
the canonical range for unsaturated loam near field capacity. The
swale 10\,cm bound is $\sim 30\%$ lower than control. The swale
40\,cm bound is two orders of magnitude smaller than control 40\,cm
and cannot be reconciled with any realistic $D_b$: it is the
signature of a regime where the diffusion mode is not the
rate-limiter at all.

\section{Why the swale 40\,cm $\tau$ blows up}

Eq.~\eqref{eq:tau-sum} says $1/\tau$ is a sum of non-negative
contributions, so $\tau \to \infty$ only when \emph{all three}
collapse simultaneously:
\begin{itemize}
\item the diffusion mode is suppressed (a saturated boundary clamps
      $\theta'$ from below, or a low-permeability layer blocks
      vertical flux);
\item gravity advection $K_b'$ vanishes (the gravitational gradient is
      balanced by a continuous lateral inflow that re-supplies water
      faster than gravity drains it); and
\item the ET sink $\lambda$ is small (deep root demand is weak at
      \SI{40}{cm}, or the available water in the perched layer is large
      compared to weekly demand).
\end{itemize}

Of these three, sustained lateral inflow from upslope mounds is the
most parsimonious single explanation, and is the one we can test
independently: it predicts a wetting-front lag at \SI{40}{cm} that
\emph{matches} \SI{10}{cm} of the same plot rather than the
\SI{40}{cm} of the control plot.

\section{What a Richards-model inversion would add}

The linear-reservoir model collapses three independent physical
processes into a single observable $\tau$. To separate them one
needs either:
\begin{enumerate}
\item A \emph{vertical-profile constraint}: two depths in the same
      column with shared $D_b$ but different surface coupling. Each
      depth provides one equation in $\{D_b, K_b', \lambda\}$. The
      current 10\,cm \& 40\,cm pair already gives this.
\item An \emph{independent ET estimate}: ATMOS\,14 vapour-pressure
      deficit and air temperature constrain $\lambda$ separately. We
      have these data on logger \texttt{Z6-19570}.
\item A \emph{forward Richards solver} (e.g.\ HYDRUS-1D / phydrus)
      with the measured rainfall as upper boundary and observed
      $\theta(z,t)$ as the inversion target, returning
      van~Genuchten / Mualem retention parameters $(\theta_r,
      \theta_s, \alpha_\text{vG}, n, K_s)$ per layer and treatment.
      Option (3) subsumes (1)--(2) and quantifies the lateral-inflow
      hypothesis as a residual.
\end{enumerate}

\end{document}
